
\documentclass[11pt,a4paper]{article}

\usepackage[margin=1in]{geometry}
\usepackage{mathpazo}
\usepackage{microtype}

\usepackage{amsmath}
\usepackage{amssymb}
\usepackage{amsthm}

\usepackage{graphicx}
\graphicspath{{../figs/}}
\usepackage{booktabs}
\usepackage{enumitem}
\usepackage[font=small,labelfont=bf,skip=4pt]{caption}
\setlength{\tabcolsep}{5pt}
\renewcommand{\arraystretch}{1.0}
\setlength{\abovedisplayskip}{5pt}
\setlength{\belowdisplayskip}{5pt}
\setlength{\abovedisplayshortskip}{3pt}
\setlength{\belowdisplayshortskip}{3pt}

\usepackage{natbib}
\usepackage[hidelinks]{hyperref}
\raggedbottom
\renewcommand{\textfraction}{0.05}
\renewcommand{\topfraction}{0.95}
\renewcommand{\bottomfraction}{0.95}
\renewcommand{\floatpagefraction}{0.85}

% ----------------------------------------------------------------------------
% Macros
% ----------------------------------------------------------------------------
\newcommand{\R}{\mathbb{R}}
\newcommand{\N}{\mathbb{N}}
\newcommand{\Z}{\mathbb{Z}}
\newcommand{\E}{\mathbb{E}}
\newcommand{\pc}{\mathbf{x}}
\newcommand{\simpcom}{\Sigma}
\newcommand{\filt}{\mathcal{F}}
\newcommand{\kpar}{\kappa}
\newcommand{\mo}{\mu}

% ----------------------------------------------------------------------------
\newtheoremstyle{definitioncolon}%
  {4pt}{4pt}{\normalfont}{}{\bfseries}{}{.5em}%
  {\thmname{#1}\thmnumber{ #2}\thmnote{: #3}}
\theoremstyle{definitioncolon}
\newtheorem{definition}{Definition}[section]
\newtheorem{claim}[definition]{Claim}

\providecommand{\pmTodo}[1]{\textbf{[TODO: #1]}}

\newtheorem*{note}{Note}
\newtheorem*{example}{Example}
\newtheorem*{remark}{Remark}

% ----------------------------------------------------------------------------
\title{Topological Feature Learning for Parameter Estimation\\
       in Neyman-Scott Processes\\[4pt]
       \large First-year progression report}
\author{Flavio Gualtieri}
\date{\today}

\begin{document}
\maketitle

% --- macros: move to preamb
% --- end macros ---

\section{Topological feature learning for Neyman--Scott process parameters}
\label{sec:paramnet}

\subsection{Motivation and problem}
\label{paramnet:sec-motivation}

Spatial point patterns -- the locations of galaxies, cells, epicentres, or
trees -- appear across the sciences, and are commonly modelled as
realizations of a spatial point process \citep{moller2004statistical}, whose
parameters are themselves interpretable scientific quantities rather than
curve-fitting artefacts. A Neyman--Scott process generates points in two
stages: unobserved \emph{parent}
locations at intensity $\kappa$, each producing a Poisson-distributed number
of \emph{offspring} (mean $\mu$) scattered around it according to a kernel
with characteristic scale $\sigma$. When such a process models galaxy
clustering, $\kappa$ is the number density of haloes hosting galaxies,
$\mu$ the mean galaxy occupancy of a halo, and $\sigma$ the spatial extent
of the galaxy distribution within a halo -- exactly the quantities a halo
occupancy model is built to constrain \pmTodo{cite halo occupancy
reference}. In forest ecology, \citet{waagepetersen2009twostep} fit Thomas
processes (a Neyman--Scott process with a Gaussian offspring kernel) to
three tropical tree species and read the cluster scale $\sigma$ directly as
a summary of seed dispersal tightness, finding that species dispersed by
exploding capsules, by wind, and by animals separate cleanly along that one
parameter: recovering $(\kappa,\mu,\sigma)$ well is how the pattern gets
turned back into an interpretable physical or ecological statement, not an
end in itself.

We consider the setting a practitioner actually faces: a single observed
point pattern, and one chance to estimate the parameters that generated it --
not repeated draws of the same process. Existing estimators for this
problem -- minimum contrast fits a discrepancy between an empirical and
theoretical summary curve \citep{diggle2003spatial}, Palm likelihood builds
a composite likelihood from the Palm intensity \citep{tanaka2008parameter},
fully Bayesian approaches sample the posterior over the latent parent
locations by MCMC \citep{moller2004statistical} -- share three limitations
that motivate this project. First, they need a closed-form summary
statistic that may not exist at all (the LGCP-Strauss process of
\citet{vihrs2022neural} has an intractable density) or may exist but fail
the regularity conditions the theory needs: \citet{waagepetersen2009twostep}
require their $K$-function to be twice continuously differentiable in the
parameters, which excludes the Mat\'ern cluster process outright. Second,
minimum contrast depends on hyperparameters ($r_{max}$, $q$, a
power-transform exponent) chosen per dataset by rule of thumb, with no
held-out criterion for whether the choice was good; \citet{waagepetersen2009twostep}
report Thomas-process estimates of $\kappa$ and $\mu$ that are ``quite
sensitive'' to $r_{max}$ and $q$, tracing this to a strong negative
correlation between $\hat\kappa$ and $\hat\mu$ under which the product
$\hat\kappa\hat\mu$ is stable while the individual estimates are not. Third,
these methods become unstable near complete spatial randomness: for a
Thomas process,
\begin{equation}
\label{paramnet:eq-kdegen}
K(t) = \pi t^2 + \frac{1-e^{-t^2/4\sigma^2}}{\kappa},
\end{equation}
so as $\sigma$ grows at fixed $t$ the second term vanishes and $K$ tends to
the Poisson curve \emph{regardless of $\kappa$} \citep{waagepetersen2009twostep}.
A discrepancy nearly flat along this trade-off yields estimates sensitive to
the optimizer's starting point; \citet{waagepetersen2009twostep} report
exactly this failure in their own simulations, weakly clustered patterns
whose estimated $K$ was indistinguishable from Poisson producing ``rather
extreme'' clustering estimates. We deliberately retain this near-Poisson
regime in our design distribution below, rather than restricting to the
well-separated-cluster case classical methods are comfortable with.

These three limitations motivate features needing no closed form and no
choice of summary statistic: a multi-scale, process-agnostic description of
the point cloud, built directly from its geometry.

\subsection{Approach}
\label{paramnet:sec-approach}

\paragraph{Persistent homology as a process-agnostic feature.}
A radially averaged statistic such as $K$ or $g$ collapses a point cloud to
a single curve indexed by scale, discarding \emph{how} density and gaps are
distributed across space. Persistent homology instead tracks a growing
family of simplicial complexes on the cloud as a scale $r$ increases: at
$r=0$ every point is its own component, and as $r$ grows, points connect,
components merge, and loops open and close. Each such feature's
\emph{birth} and \emph{death} scale is recorded as a point $(b,d)$ in a
\emph{persistence diagram} (points with $d\!-\!b$ small are usually treated
as noise). The $0$-dimensional diagram records when clusters merge; the
$1$-dimensional diagram records voids and gaps that persist across scales,
capturing empty or repulsive structure a radially averaged statistic sees
only indirectly. Neither needs a closed form or an assumption about the
generating process family. A diagram is a variable-size multiset rather
than a vector, so it must be \emph{vectorized} before it can feed a neural
network; we use \emph{persistence images} \citep{adams2017persistence},
which smooth each diagram point with a Gaussian kernel and integrate onto a
fixed pixel grid. Betti curves \citep{umeda2017time}, persistence
landscapes \citep{bubenik2015landscapes}, and persistence silhouettes
\citep{chazal2014silhouette} are alternative vectorizations we compare
against directly (Figure~\ref{paramnet:fig-vectorizations}); none improved
on persistence images under our architecture and design distribution
(Section~\ref{paramnet:sec-results}).

\begin{figure}[tbp]
    \centering
    \includegraphics[width=\textwidth]{vectorization_schematic.pdf}
    \caption{One point cloud $\rightarrow$ its persistence diagrams
    $\rightarrow$ the vectorizations compared: persistence-image raster,
    landscape stack, silhouette curve, Betti curve, and a persistence-statistics
    summary. Top: a tight-cluster design ($\kappa=50,\mu=4,\sigma=0.01$);
    bottom: a near-Poisson design ($\kappa=4,\mu=5,\sigma=0.5$) -- the
    regime Section~\ref{paramnet:sec-motivation} argues classical methods
    handle badly.}
    \label{paramnet:fig-vectorizations}
\end{figure}

\paragraph{Simulator and design distribution.}
Training data are pairs $(\mathbf{x}_i,\theta_i)$, each a single
realization: parents form a Poisson process on a slightly expanded window
(so clusters whose parent lies just outside are not truncated), each parent
independently produces a Poisson-distributed number of offspring, and each
offspring is displaced from its parent by a draw from the kernel --
Gaussian for Thomas, and, for nested Thomas, a second Thomas process
standing in for the parent layer itself. Sampling $\theta_i$ uniformly over
$(\kappa,\mu,\sigma)$ is inadequate because these processes separate into
qualitative regimes occupying curved regions of that space, wasting most of
the sampled mass on configurations that are near-Poisson or degenerate. We
instead reparameterize into coordinates carrying the same meaning for every
Neyman--Scott family member: the expected offspring intensity
$\lambda=\kappa\mu$, the parent intensity $\kappa$, and a dimensionless
overlap index
\begin{equation}
\label{paramnet:eq-reparam}
\nu = \sigma\sqrt{\kappa},
\end{equation}
the ratio of cluster radius to mean inter-parent spacing, and sample
log-uniformly in $(\kappa,\lambda,\nu)$ subject to a point budget
($\lambda\in[150,800]$, since persistence computation scales badly with
cloud size), a resolvability constraint ($\kappa\in[15,120]$), and an
overlap range $\nu\in[0.10,0.90]$ whose upper end deliberately keeps the
near-Poisson regime of Section~\ref{paramnet:sec-motivation} in training
rather than excluding it.

\paragraph{ParamNet.}
The model maps a cloud to a parameter estimate by
\begin{equation}
\label{paramnet:eq-pipeline}
\text{cloud} \;\xrightarrow{\;\text{vectorize}\;}\;
\text{persistence image} \;\xrightarrow{\;\text{encoder}\;}\;
\text{embedding} \;\xrightarrow{\;\text{FCNN}\;}\;
\hat\theta,
\end{equation}
following the neural-estimator pattern introduced by
\citet{vihrs2022neural}; our contribution is the feature-extraction and
calibration stage, its evaluation on a genuinely two-level process (nested
Thomas), and a direct comparison of vectorizations. Images come from a
density-sensitive DTM$_k$ filtration rather than the more common
Vietoris--Rips filtration: Rips uses raw inter-point distances and is
sensitive only to the sparsest connection between two regions, whereas
DTM$_k$ uses each point's average distance to its $k$ nearest neighbours
and so responds directly to local density -- exactly what a Neyman--Scott
process's parameters govern \citep{anai2020dtm}. The Thomas model uses a
single scale ($k=5$); nested Thomas fuses three scales ($k\in\{5,10,15\}$)
through one shared convolutional encoder, which works better than giving
each scale its own encoder (Section~\ref{paramnet:sec-results}). Before
rasterizing, each diagram's axes are calibrated to a shared range fit on
the training diagrams only -- a per-diagram coverage-quantile rule, so one
outlier diagram cannot compress every other diagram's range toward the
origin -- then frozen before touching validation or test data. The encoder
is a small CoordConv
\citep{liu2018coordconv} convolutional stack: CoordConv appends each
pixel's normalized coordinates as extra input channels, giving the
otherwise translation-equivariant convolutions access to absolute position,
which matters because a persistence image's axes are not arbitrary --
birth encodes local density, persistence encodes separation -- so a
pixel's location is informative in its own right. The resulting embedding
is concatenated with $\log N(\mathbf{x})$ and passed to a small
feed-forward head producing $\hat\theta$. The model is trained by
minimizing mean squared error between $\hat\theta$ and $\theta$ in
log-normalized parameter space (each parameter log-transformed and
$z$-scored on the training split), so parameters spanning different orders
of magnitude contribute comparably to the loss. Data are split by
parameter vector, so no $\theta$ value appears in both training and
evaluation, and error is reported in those same log-normalized units.

\subsection{Results}
\label{paramnet:sec-results}

We compare ParamNet against minimum contrast fit separately to $K$ and to
$g$ \citep{diggle2003spatial}, and against \texttt{vihrs}, our
re-implementation of the neural estimator of \citet{vihrs2022neural}: the
same pipeline with the persistence-image feature replaced by a 1-D CNN over
the classical, edge-corrected $L(r)-r$ curve plus $\log N$. A label-shuffle
control -- the same model retrained with training labels permuted -- lands
at $\mathcal{L}\approx1.0$ on both processes and serves as a chance-level
ceiling in Table~\ref{paramnet:tab-comparison}.

\begin{table}[tbp]
\centering
\small
\begin{tabular}{lcc}
\toprule
Method & Thomas $\mathcal{L}$ & Nested-Thomas $\mathcal{L}$ \\
\midrule
ParamNet                & $0.125 \pm 0.007$ & $0.192 \pm 0.010$ \\
\texttt{vihrs}           & $0.115 \pm 0.007$ & $0.217 \pm 0.006$ \\
Minimum contrast, $K$    & $0.855 \pm 0.682$ & ---$^{a}$ \\
Minimum contrast, $g$    & $6.199 \pm 0.339$ & ---$^{a}$ \\
\midrule
\emph{Chance level (label shuffle)} & \emph{$1.003 \pm 0.019$} & \emph{$1.006 \pm 0.020$} \\
\bottomrule
\end{tabular}
\caption{Mean squared error in log-normalized parameter space, $n=10$ seeds,
shared test split per process. $^{a}$A closed-form $K$/$g$ for nested Thomas
was derived and fitted, but the derivation has not yet been independently
validated, so we do not report the resulting numbers as a finding.}
\label{paramnet:tab-comparison}
\end{table}

Persistence images carry real signal about Neyman--Scott parameters: both
processes land roughly an order of magnitude below the label-shuffle
ceiling, and that signal sits in the images rather than in the point count
-- training on $\log N$ alone collapses to the same $\approx\!0.9$--$1.0$
chance level as label-shuffling, while removing $\log N$ from the full
model costs little on nested Thomas ($0.192\rightarrow0.203$) and is merely
noisy on Thomas. We report the comparison honestly rather than as an
unqualified win, though: \texttt{vihrs} beats ParamNet on every Thomas
parameter ($\mathcal{L}^\kappa=0.192$ vs.\ $0.201$, $\mathcal{L}^\mu=0.122$
vs.\ $0.134$, $\mathcal{L}^\sigma=0.029$ vs.\ $0.040$), so a hand-built
radial summary beats our topological pipeline outright on the single-level
process. On nested Thomas the picture inverts only partially: ParamNet
gains an advantage on the density and count targets (parent intensity,
meta-offspring, mean offspring), while \texttt{vihrs} keeps its edge on
both scale parameters -- $L(r)-r$ is built to resolve a radius directly, a
density-based topological feature is not, and it is only on the more
complex two-level process that the trade-off favours our pipeline. A
second, structural pattern holds regardless of estimator or process: parent
intensity $\kappa$ is the hardest target by a wide margin and cluster scale
the easiest, reading as a property of the estimation problem rather than of
any one feature set.

Both neural methods dominate minimum contrast on Thomas by roughly an order
of magnitude, but for different reasons. On $K$, six of ten seeds land at
$0.365$--$0.399$ (close to \texttt{vihrs}) while the other four land at
$0.903$--$2.056$ -- seed-to-seed instability in the near-Poisson regime of
Section~\ref{paramnet:sec-motivation}, not a uniformly bad fit. On $g$,
every seed is worse than chance, pointing to a systematic bias rather than
scattered non-convergence: our contrast exponent ($c=1$, chosen by analogy
with the $K$-based rule of thumb) was never tuned against the $c=1/4$ value
validated for aggregated patterns \citep{diggle2003spatial}; we do not read
it as a general verdict on contrast estimation. The same comparison was
attempted for nested Thomas this month, against a closed-form $K$/$g$
derived for the two-level process, but the derivation has not been checked
independently and the resulting fits are far outside a plausible range, so
we treat this baseline as still outstanding rather than report numbers.

We also compared three alternative vectorizations (Betti curves,
landscapes, silhouettes) against persistence images under an identical
encoder and design distribution; none improved on either process. Fusing
the three DTM scales through one shared encoder likewise beats giving each
scale its own encoder ($0.192$ vs.\ $0.222$ nested-Thomas $\mathcal{L}$),
which is why ParamNet shares an encoder rather than using independent
branches.

\subsection{Limitations and future directions}
\label{paramnet:sec-future}

\paragraph{Limitations.}
\begin{itemize}[leftmargin=*,itemsep=1pt,topsep=2pt]
\item \textbf{No advantage on the single-level process.} \texttt{vihrs}
    wins on every Thomas parameter and on both scale parameters of nested
    Thomas; our advantage is confined to the two-level process's density
    and count targets, and to three alternative vectorizations under one
    architecture and design distribution -- evidence about this setting,
    not a general claim.
\item \textbf{Multiple models, and expensive ones.} The number of
    filtration scales fused is a per-process choice made by comparing
    validation losses, undercutting the process-agnostic framing behind
    this work; always fusing three scales costs three times the
    persistence-computation cost for a Thomas change within seed noise, and
    persistence computation already scales badly with point count, which is
    why the design distribution caps the expected count at 800.
\item \textbf{No irreducible error floor yet.} The natural lower reference
    for every number above -- the dispersion of the best possible estimator
    across repeated realizations of the same $\theta$ -- has been attempted
    but not obtained: the estimation script ran to completion but returned
    \texttt{NaN} throughout, traced to a suspected label-indexing bug not
    yet fixed.
\end{itemize}

\paragraph{Future directions.}
\begin{itemize}[leftmargin=*,itemsep=1pt,topsep=2pt]
\item A universal, process-agnostic encoder, rather than a per-process
    filtration-scale choice; a small three-seed comparison also suggests a
    multiparameter persistence feature (a joint density--radius
    bifiltration fused with the \texttt{vihrs} branch) may close part of
    the Thomas gap, though not yet validated at the scale above.
\item Inhomogeneous and non-stationary processes, where classical closed
    forms are least available and a topological approach has the most to
    gain: the natural comparator, the two-step procedure of
    \citet{waagepetersen2009twostep}, needs observed covariates and a
    correctly specified intensity, whereas a simulation-trained estimator
    needs neither.
\item Real data: galaxy catalogues, motivated by
    Section~\ref{paramnet:sec-motivation}, and the Barro Colorado Island
    tree data of \citet{waagepetersen2009twostep}, already fitted with a
    Thomas process.
\end{itemize}


\end{document}